\subsection{JADE agents}
\label{sec:jade-basics}
\slidesrc{JADE: 4--9}

\begin{nbdef}{JADE (slides 4--5)}
\textbf{JADE} (Java Agent DEvelopment Framework; Bellifemine, Caire, Greenwood,
\emph{Developing Multi-Agent Systems with JADE}, Wiley 2007; latest release December 2022)
is a software framework fully implemented in Java. It simplifies the implementation of MAS
through
\begin{itemize}[nosep]
  \item a \textbf{middleware} that complies with the FIPA specifications, and
  \item a set of \textbf{tools} that support debugging and deployment (e.g.\ the sniffer
        that draws message exchanges, Section~\ref{sec:uses-of-mas}).
\end{itemize}
The agent platform can be \textbf{distributed} across machines, which need not share the
same operating system, and the configuration can be controlled via a remote GUI. It can be
changed at run time by \textbf{moving agents} from one machine to another, as and when
required.
\end{nbdef}

\begin{center}
\begin{tikzpicture}[font=\small, >=Stealth,
  host/.style={draw, thick, rounded corners, minimum width=4.3cm, minimum height=2.0cm, fill=#1},
  ag/.style={circle, draw, thick, fill=nbgreen!25, minimum size=0.6cm, font=\scriptsize}]
  \node[host=nbblue!6] (h1) at (0,0) {};
  \node[font=\scriptsize\bfseries, anchor=north] at (h1.north) {Main container (Linux)};
  \node[ag] at (-1.0,-0.2) {A1}; \node[ag] at (0,-0.2) {AMS}; \node[ag] at (1.0,-0.2) {DF};
  \node[host=nbblue!6] (h2) at (4.8,0) {};
  \node[font=\scriptsize\bfseries, anchor=north] at (h2.north) {Container 1 (Windows)};
  \node[ag] (a2) at (4.3,-0.2) {A2}; \node[ag] at (5.3,-0.2) {A3};
  \node[host=nbblue!6] (h3) at (9.6,0) {};
  \node[font=\scriptsize\bfseries, anchor=north] at (h3.north) {Container 2 (macOS)};
  \node[ag] (a4) at (9.6,-0.2) {A4};
  \draw[very thick, nbred!70] (-2.2,-1.6) -- (11.8,-1.6) node[midway, below, font=\scriptsize, text=black]{one JADE platform = distributed runtime (FIPA message transport)};
  \foreach \h in {h1,h2,h3} \draw[thick, nbred!70] (\h.south) -- (\h.south |- 0,-1.6);
  \draw[->, thick, dashed, nborange] (a2.south) to[bend right=20] node[pos=0.35, above, font=\scriptsize]{agent mobility} (a4.south);
\end{tikzpicture}
\captionof{figure}{A JADE platform spread over three machines (added drawing). The main
container hosts the AMS (white pages) and the DF (yellow pages); agents can move between
containers at run time.}
\end{center}

\subsubsection{The HelloWorld agent and AIDs}
\slidesrc{JADE: 6--7}

A \textbf{type of agent} is created by extending the \code{jade.core.Agent} class and
redefining the \code{setup()} method (called when the agent starts). Each agent instance is
identified by an \textbf{AID} (\code{jade.core.AID}), composed of a unique \emph{name} plus
some \emph{addresses}; an agent retrieves its own AID with \code{getAID()}.
\begin{minted}{java}
import jade.core.Agent;

public class HalloWorldAgent extends Agent {
    protected void setup() {
        System.out.println("Hallo World! my name is " + getAID().getName());
    }
}
\end{minted}

\begin{nbfact}{Names, GUID and addresses (slide 7)}
\begin{itemize}[nosep]
  \item Agent names have the form \code{<local-name>@<platform-name>}, and the complete
        name must be \textbf{globally unique} (GUID).
  \item The default platform name is \code{<main-host>:<main-port>/JADE}; it can be set with
        the \code{-name} option.
  \item Within a single platform, agents are referred to by their (local) names only.
  \item Given a name, an AID is created as \code{new AID(localname, AID.ISLOCALNAME)} or
        \code{new AID(name, AID.ISGUID)}.
  \item The addresses in an AID are those of the platform's MTPs (message transport
        protocols) and are used \textbf{only} for communication between agents on
        \emph{different} FIPA platforms.
\end{itemize}
Example: an agent started as \code{peter} on the platform \code{myhost:1099/JADE} has GUID
\code{peter@myhost:1099/JADE}; \code{new AID("peter", AID.ISLOCALNAME)} refers to it from
the same platform.
\end{nbfact}

\subsubsection{Arguments and termination}
\slidesrc{JADE: 8--9}

It is possible to pass arguments to an agent at start-up,
\code{java jade.Boot ... a:myPackage.MyAgent(arg1 arg2)}, and retrieve them with
\code{getArguments()}. An agent terminates when its \code{doDelete()} method is called; on
termination its \code{takeDown()} method is invoked (intended for clean-up operations).
\begin{minted}{java}
protected void setup() {
    System.out.println("Hallo World! my name is " + getAID().getName());
    Object[] args = getArguments();
    if (args != null) {
        System.out.println("My arguments are:");
        for (int i = 0; i < args.length; ++i) {
            System.out.println("- " + args[i]);
        }
    }
    doDelete();        // terminate immediately after greeting
}

protected void takeDown() {
    System.out.println("Bye...");
}
\end{minted}

\begin{nbexample}{what is printed}
\notinslides
Started as \code{a:myPackage.MyAgent(red 42)} on platform \code{myhost:1099/JADE}:
\begin{minted}[linenos=false]{text}
Hallo World! my name is a@myhost:1099/JADE
My arguments are:
- red
- 42
Bye...
\end{minted}
\code{doDelete()} does not stop the agent in the middle of \code{setup()}: \code{setup()}
completes, then the agent sees the deletion request, calls \code{takeDown()} and dies (see
the execution model in the next sub-section). Arguments arrive as \code{Object}s (strings
from the command line).
\end{nbexample}
